\begin{table}[H]
\centering
\caption{Fama--French (2010) Bootstrap: Actual vs.\ Simulated $t(\hat{\alpha})$ Percentiles -- FF Subperiod}
\label{tab:bootstrap_tails_FF}
\begin{threeparttable}
\begin{tabular}{rrrrl}
\toprule
Percentile & Actual $t(\hat{\alpha})$ & Simulated Mean & Prob.\ Luck & Interpretation \\
\midrule
 1 & -3.297 & -2.246 & 1.4\% & Worse than luck (significant) \\
 5 & -2.437 & -1.463 & 0.1\% & Worse than luck (significant) \\
 10 & -1.994 & -1.111 & 0.1\% & Worse than luck (significant) \\
 50 & -0.472 & 0.014 & 1.4\% & Indistinguishable from luck \\
 90 & 1.106 & 1.129 & 48.6\% & Consistent with zero-skill \\
 95 & 1.559 & 1.481 & 63.5\% & Consistent with zero-skill \\
 99 & 2.370 & 2.296 & 61.6\% & Consistent with zero-skill \\
\bottomrule
\end{tabular}
\begin{tablenotes}
\item Bootstrap procedure follows \textcite{FamaFrench2010}. Sample: actively managed funds, Jan 1995--Sep 2006 (Fama--French 2010 subperiod), performance-comparison subsample per flagged\_funds.xlsx, minimum 24 monthly observations ($N = 805$ funds). For each fund, estimated monthly alpha is subtracted from the excess return series to construct a zero-alpha null return. In each of $B = 10{,}000$ bootstrap iterations, calendar months are resampled with replacement, preserving cross-sectional factor return dependence. The \textcite{Carhart1997} four-factor model is re-estimated on each resampled series. \textit{Actual} $t(\hat{\alpha})$: percentile of the empirical $t$-statistic distribution across active funds. \textit{Simulated Mean}: average of that percentile across all iterations. \textit{Prob.\ Luck}: fraction of iterations in which the simulated percentile falls below the actual value; values below 5\% at lower percentiles indicate underperformance unlikely to be explained by luck alone. Newey--West standard errors with a 6-month lag are used throughout. Compare with \textcite{FamaFrench2010}, Table~III.
\end{tablenotes}
\end{threeparttable}
\end{table}
